% Part: history
% Chapter: biographies 
% Section: alan-turing 

\documentclass[../../../include/open-logic-section]{subfiles}

\begin{document}

\olfileid[hi]{his}{bio}{tur} 

\olsection{एलन ट्यूरिंग}

एलन ट्यूरिंग का जन्म 23 जून 1912 को लंदन के मैडा वेल में हुआ। उन्हें
सैद्धांतिक कंप्यूटर विज्ञान का जनक माना जाता है। भौतिक विज्ञानों और गणित में
ट्यूरिंग की रुचि कम आयु में ही आरंभ हो गई थी। किंतु बाल्यावस्था में उनके
विद्यालयों में उनकी रुचियों को समुचित स्थान नहीं मिला; वहाँ साहित्य और
शास्त्रीय विषयों पर बल दिया जाता था। फलतः विद्यालय में उनका प्रदर्शन खराब
रहा और उनके कई शिक्षकों ने उन्हें फटकारा।

\olphoto{turing-alan}{एलन ट्यूरिंग}

ट्यूरिंग ने स्नातक विद्यार्थी के रूप में किंग्स कॉलेज, कैम्ब्रिज में गणित
पढ़ा। 1936 में उन्होंने गणनीय फलन की संकल्पना को ठीक-ठीक परिभाषित करने और
निर्णय समस्या की अनिर्णेयता सिद्ध करने के प्रयास में वह यंत्र विकसित किया
जिसे अब ट्यूरिंग मशीन कहते हैं। अलॉन्ज़ो चर्च अपने लैम्ब्डा कलन द्वारा इस
परिणाम को उनसे पहले सिद्ध कर चुके थे। फिर भी चर्च के परिणाम के संदर्भ सहित
ट्यूरिंग का शोधपत्र प्रकाशित हुआ। चर्च ने ट्यूरिंग को प्रिंसटन आमंत्रित किया,
जहाँ वे 1936--1938 तक रहे और चर्च के निर्देशन में डॉक्टरेट प्राप्त की।

तर्कशास्त्र में रुचि के बावजूद भौतिक विज्ञानों में ट्यूरिंग की आरंभिक रुचि
प्रबल बनी रही। द्वितीय विश्व युद्ध में ब्लेचली पार्क स्थित ब्रिटिश
कूटविश्लेषण विभाग में सेवा के दौरान उनके व्यावहारिक कौशल का उपयोग हुआ। जर्मन
नौसैनिक संचार में प्रयुक्त एनिग्मा कूट को तोड़ने में ट्यूरिंग की केंद्रीय
भूमिका थी। इलेक्ट्रॉनिक मशीनों के आगमन के साथ सांख्यिकी और कूटलेखन में उनकी
विशेषज्ञता ने दल को ``बॉम्ब'' नामक विकूटन मशीन बनाकर कूट तोड़ने में सक्षम
किया। उनके विचारों ने विश्व के पहले प्रोग्राम-योग्य इलेक्ट्रॉनिक कंप्यूटर,
कोलॉसस, के निर्माण में भी सहायता की; इसका उपयोग ब्लेचली पार्क में जर्मन
लोरेंज़ कूट तोड़ने के लिए भी हुआ।

ट्यूरिंग समलैंगिक थे। फिर भी 1942 में उन्होंने ब्लेचली पार्क की अपनी
सहकर्मी जोन क्लार्क को विवाह का प्रस्ताव दिया, किंतु बाद में सगाई तोड़कर
उन्हें बताया कि वे समलैंगिक हैं। उनके जीवन में कई प्रेम-संबंध रहे, यद्यपि उस
समय ब्रिटेन में समलैंगिक कृत्य दंडनीय अपराध थे। 1952 में ट्यूरिंग के तत्कालीन
प्रेमी के एक मित्र ने उनके घर में चोरी की। पुलिस में शिकायत करते समय ट्यूरिंग
ने अपना समलैंगिक संबंध स्वीकार किया, क्योंकि उन्हें लगा कि सरकार समलैंगिक
कृत्यों को वैध बनाने वाली है। यह सच नहीं था और उन पर घोर अभद्रता का आरोप
लगा। कारावास के स्थान पर ट्यूरिंग ने कामेच्छा घटाने वाला हार्मोन उपचार चुना।
8 जून 1954 को वे सायनाइड की अत्यधिक मात्रा से मृत पाए गए---सबसे संभवतः यह
आत्महत्या थी। 2013 में महारानी एलिज़ाबेथ~द्वितीय ने उन्हें शाही क्षमादान
दिया।

\begin{reading}
एलन ट्यूरिंग की विस्तृत जीवनी के लिए \citet{Hodges2014} देखें। ट्यूरिंग के
जीवन और कार्य से प्रेरित नाटक \emph{Breaking the Code} का 1996 में टीवी
निर्माण हुआ, जिसमें डेरेक जैकोबी ने ट्यूरिंग की भूमिका निभाई। बेनेडिक्ट
कम्बरबैच और कीरा नाइटली अभिनीत अकादमी पुरस्कार-नामांकित फिल्म \emph{The
Imitation Game} भी एलन ट्यूरिंग के जीवन और ब्लेचली पार्क में उनके समय पर
शिथिल रूप से आधारित है \citep{Imitation2014}।

\citet{Radiolab2012} में ट्यूरिंग के जीवन और कार्य पर कई पॉडकास्ट हैं। बीबीसी
होराइज़न का वृत्तचित्र \emph{The Strange Life and Death of Dr.~Turing}
ऑनलाइन देखने को उपलब्ध है \citep{Sykes1992}। \citet{Theelen2012} एक चलती
हुई लेगो ट्यूरिंग मशीन का लघु वीडियो है---इसे 2012 में ट्यूरिंग की जन्मशती
के सम्मान में बनाया गया था।

ट्यूरिंग मशीनों और निर्णय समस्या पर ट्यूरिंग का मूल शोधपत्र
\citet{Turing1937} है।
\end{reading}

\end{document}
